\documentclass[11pt,a4paper]{article}
\usepackage[T1]{fontenc}
\usepackage[utf8]{inputenc}
\usepackage{lmodern,amsmath,amssymb}
\usepackage[margin=25mm]{geometry}
\usepackage[hidelinks]{hyperref}
\hypersetup{pdftitle={Reciprocal exchange and mechanical normalization},pdfauthor={navstokgap research working notes}}
\newcommand{\tightlist}{\setlength{\itemsep}{0pt}\setlength{\parskip}{0pt}}
\setlength{\emergencystretch}{3em}
\setcounter{secnumdepth}{0}
\begin{document}
\hypertarget{reciprocal-exchange-fixes-relative-mechanical-scales}{%
\section{Reciprocal exchange fixes relative mechanical
scales}\label{reciprocal-exchange-fixes-relative-mechanical-scales}}

Two measured acceleration responses determine an inertia ratio when both
arise from a shared conservative spring. A connected network determines
all such ratios if its cycle products are consistent. One common
multiplier of inertias, spring constants, energies and actions remains
free. Reciprocal exchange therefore supplies a relative-normalization
mechanism, with an explicit residual absolute-scale freedom.

Status: Q06 exploratory derivation, not an accepted-ledger claim. The
model and written consistency review below require independent proof
review and a bounded literature comparison before promotion. No novelty
claim is made.

\hypertarget{specify-the-response-before-choosing-masses}{%
\subsection{1. Specify the response before choosing
masses}\label{specify-the-response-before-choosing-masses}}

Let two coordinates \(x,y\in\mathbb R\) have the same length units and
share physical time \(t\). Suppose their measured equations are

\[\ddot x=-\alpha(x-y),\qquad
\ddot y=\beta(x-y),\qquad \alpha,\beta>0. \tag{1}\]

Both coefficients have units time\(^{-2}\). Isolated free motion would
leave the sectors' positive inertias \(m_1,m_2\) independently
arbitrary. Now require the specific mechanical completion

\[L=\tfrac12m_1\dot x^2+\tfrac12m_2\dot y^2
-\tfrac12k(x-y)^2,\qquad k>0. \tag{2}\]

Here masses have mass units and \(k\) has mass/time\(^2\) units. This
assumes a diagonal constant kinetic energy and one reciprocal
position-dependent interaction. It is a physical model premise, not a
consequence of observing arbitrary coupled motion. Its forces are
\(F_1=-k(x-y)\) and \(F_2=k(x-y)\). The Euler--Lagrange equations agree
with (1) exactly when

\[m_1\alpha=m_2\beta=k,
\qquad \frac{m_2}{m_1}=\frac\alpha\beta. \tag{3}\]

Thus fixed responses admit precisely the positive family
\((m_1,m_2,k)=m_1(1,\alpha/\beta,\alpha)\), with the entries interpreted
in their stated units. Coupling constrains relative normalization; equal
inertias follow only if the two acceleration responses are equal.

\hypertarget{energy-exchange-and-the-remaining-action-freedom}{%
\subsection{2. Energy exchange and the remaining action
freedom}\label{energy-exchange-and-the-remaining-action-freedom}}

Writing \(r=x-y\), the conserved energy and total momentum are

\[H=\tfrac12m_1\dot x^2+\tfrac12m_2\dot y^2+\tfrac12kr^2,
\qquad P=m_1\dot x+m_2\dot y. \tag{4}\]

Indeed \(\dot T_1=-kr\dot x\), \(\dot T_2=kr\dot y\), and
\(\dot V=kr(\dot x-\dot y)\), so \(\dot H=0\); the two forces give
\(\dot P=0\). To display an exchange current, assign half the
interaction energy to each sector: \(E_1=T_1+V/2\), \(E_2=T_2+V/2\).
Then

\[\dot E_1=-\tfrac{k}{2}r(\dot x+\dot y)=-\dot E_2. \tag{5}\]

The half split is an accounting convention; conservation of \(H\) is
independent of it. Generic initial data give nonzero exchange. Measuring
power in an independently calibrated energy unit would add information
beyond (1).

For every dimensionless \(\eta>0\), replace \((m_1,m_2,k)\) by
\((\eta m_1,\eta m_2,\eta k)\). Equations (1), all coordinate histories
with the same initial positions and velocities, and their periods remain
unchanged. \(L,H,P,E_i\) and the currents in (5) scale by \(\eta\). So
does the trajectory action \(S=\int L\,dt\) on any fixed time interval.
Independent rescaling of only one sector generally violates (3); common
rescaling does not.

An on-shell positive action diagnostic makes this freedom visible
without cancellation in \(S\). Set \(M=m_1+m_2\), \(\mu=m_1m_2/M\) and
\(X=(m_1x+m_2y)/M\). Then

\[L=\tfrac12M\dot X^2+\tfrac12\mu\dot r^2-\tfrac12kr^2,
\qquad \omega^2=k/\mu=\alpha+\beta. \tag{6}\]

In the stationary centre frame choose \(r=A\cos(\omega t)\), \(A>0\).
With \(p_r=\mu\dot r\), its closed-orbit action is

\[I=\frac1{2\pi}\oint p_r\,dr
=\frac12\mu\omega A^2=\frac{E_{\rm rel}}\omega. \tag{7}\]

The integral follows from \(p_r\dot r=\mu A^2\omega^2\sin^2(\omega t)\)
over one period. It has action units. Common rescaling sends \(I\) to
\(\eta I\) at fixed trajectory. Even for fixed masses and spring, the
admitted small amplitudes \(A\downarrow0\) send \(I\) to zero. The full
Lagrangian action over this stationary-centre period is zero because the
kinetic and potential integrals coincide; (7) must not be identified
with that signed integral.

The relative frequency is strictly positive for this fixed two-body
model. The full system also has a free centre coordinate, so it has a
zero-frequency mode. This is a finite-dimensional classical frequency
statement, with no volume or continuum claim and no quantum energy gap.

\hypertarget{a-network-consistency-test}{%
\subsection{3. A network consistency
test}\label{a-network-consistency-test}}

The relative mechanism extends to a finite connected undirected graph of
scalar coordinates, all measured in the same length and time units. On
every edge \(\{i,j\}\) suppose both directed responses
\(a_{ij},a_{ji}>0\) are known:

\[\ddot q_i=-\sum_{j\sim i}a_{ij}(q_i-q_j). \tag{8}\]

Seek only the mechanical class

\[L=\tfrac12\sum_i m_i\dot q_i^2
-\tfrac12\sum_{\{i,j\}}k_{ij}(q_i-q_j)^2,
\qquad m_i>0,\quad k_{ij}=k_{ji}>0. \tag{9}\]

Comparison of coefficients requires \(m_i a_{ij}=m_j a_{ji}=k_{ij}\).
Hence the necessary cycle condition is

\[\prod_{\ell=0}^{s-1}
\frac{a_{i_\ell i_{\ell+1}}}{a_{i_{\ell+1}i_\ell}}=1,
\qquad i_s=i_0. \tag{10}\]

It is sufficient as well: choose any \(m_0>0\), transport masses along a
spanning tree using \(m_j/m_i=a_{ij}/a_{ji}\), and apply (10) on each
remaining edge's fundamental cycle. This makes its transported ratio
agree with the edge ratio. All masses are positive and define the
required springs. Any two solutions have the same mass ratio along every
edge; connectivity makes their quotient one common constant. A
disconnected graph instead retains one constant per connected component.

For a concrete failed test take a triangle with \(a_{12}=2\gamma\) and
all other directed edge coefficients \(\gamma>0\). The cycle
\(1\to2\to3\to1\) has product 2. The required ratios would give
\(m_2=2m_1\), \(m_3=m_2\), and \(m_1=m_3\), a contradiction. Such
responses cannot have completion (9). This rejects that diagonal-inertia
reciprocal-spring model; it does not exclude completions with additional
variables or a different kinetic form.

\hypertarget{decision-and-written-consistency-review}{%
\subsection{4. Decision and written consistency
review}\label{decision-and-written-consistency-review}}

The positive outcome is a concrete relative calibration mechanism:
reciprocal exchange joins formerly independent normalizations. Its cycle
condition is a falsifiable constraint on measured responses. It does not
dynamically drive masses toward a ratio; it identifies which fixed
masses are compatible with the assumed energy and observed equations. A
calibrated reference mass, spring force or energy transfer would fix the
common multiplier empirically. It would supply the missing normalization
as an input.

The coordinator's written checks are (3) by direct variation, (4)--(5)
by the displayed work balance, (6)--(7) by the centre/relative change of
coordinates and period integral, and (10) by path consistency. Scaling
changes mechanical parameters rather than coordinate units: a fixed
external force or mass standard would distinguish the models. Positive
\(\eta\) and positive amplitudes approach zero action without using a
degenerate mass endpoint. No numerical or symbolic verification scripts
were used.

Q05 left open whether mutual interaction could remove its common action
coefficient. Within (2) and (9), fixed mutual responses remove relative
freedoms but retain the common one. Park further spring-network variants
unless a concrete measured response or proposed interaction violates
this completion and changes a named physical premise. For the quantum
track, the next distinct test is global spin-action patching:
distinguish classical symplectic consistency from consistency of a
proposed phase \(\exp(iS/K)\), and locate which premise, if any, forces
integrality or fixes \(K\).

Source context is the \href{stabilized-topology-action-scale.md}{Q05
normalization calculation} and
\href{composition-universality.md}{conditional composition result}.
Their source audits are retained for their own claims. The present
elementary mechanical construction has no new prior-art comparison yet;
this exploratory note is not a ledger promotion.
\end{document}
